% Research notes: time, recurrent state and the serving stack (moved unchanged from the paper's
% Appendix E, except cross-references; owner: appendix editor).
\section{Time, recurrent state and the serving stack}\label{notes:stack}
Appendix~\paperref{app:stack} of the paper measures where the time of the stock engine's decode step and speculative cycle goes and defines the serving protocol behind its frontier. This section holds the rest of that appendix: the yardstick for every lever, the ranked costs, derived ceilings for stacks of levers, how speculation interacts with the recurrent state, the proposals P4 and P5, the levers measured so far and the two stacks. Every measurement here is of the stock engine or of a lever applied to it; none uses the certified head.

\subsection{Useful progress per unit time}\label{notes:rate}
The yardstick for every lever is useful progress per unit time.\label{sec:rate} Let a cycle commit $R$ output tokens in elapsed time $T>0$. Under a stationary regenerative model with finite expectations the long-run rate is
\begin{equation}\lambda=\frac{\E R}{\E T},\label{eq:rate}\end{equation}
the ratio of accumulated reward to accumulated time, not $\E[R/T]$, and not a queueing-latency model. A change improves the rate exactly when $\E R_1/\E R_0>\E T_1/\E T_0$. Higher acceptance does not compensate for arbitrary overhead, and lower acceptance can be worthwhile when the saved work exceeds the lost progress. The denominator includes drafting, verification, the head, refinement, completion, state commit, non-overlapped dispatch and control. A fixed-width target pass evaluates all admitted positions even when the first proposal is rejected, and overlapped kernels require a critical-path measurement rather than a sum of kernel durations. If positive intervals enclosed the expected progress and time of a candidate and a baseline, $R_1^-/T_1^+>R_0^+/T_0^-$ (with $T_0^->0$) would certify $\lambda_1>\lambda_0$; benchmark confidence intervals are not deterministic enclosures, so the inequality says what must be estimated rather than certifying a speedup.

\subsection{Ranked costs}\label{notes:ranked}
Table~\ref{tab:profile-ranked} ranks the measured costs of the stock engine (Appendix~\paperref{app:profile} of the paper) by the Amdahl ceiling $1/(1-f)$ of removing each one, a bound for removing a component under otherwise unchanged execution, not a predicted gain. The other layers of the stack leave more headroom than the head (H6).

\paragraph{The plain step beyond the head} If every weight GEMM ran at the head GEMM's rate the step would be 469~\textmu s ($B=1$) to 923~\textmu s ($B=128$) shorter~\ev{profile-attr}. The recurrent kernel's time is linear in the batch, 28.9~\textmu s per request per step, which implies 3.47--3.48~TB/s at $B=32$ and 128: it streams the state once each way at 92\% of the read peak. In bytes the state equals the weights at $B=84$; in time the crossover comes later, near $B=110$--$120$, because the weight GEMMs run below the peak while the state kernel runs near it, and at $B=128$ the recurrent kernel alone (3.70~ms of wall time) already takes longer than all weight GEMMs together, head included (3.47~ms)~\ev{profile-attr,profile-bytes}. The host is not on the critical path: each graph launch is issued 1.7~ms ($B=1$) to 5.9~ms ($B=128$) before the GPU starts it, and the GPU is busy 95.7--98.3\% of the step. The margin is finite, though. At $B=1$ the scheduler thread spends only 0.73~ms of a 3.54~ms step blocked on the GPU, so a GPU-side speedup much above $1.25\times$ at batch one would make the host the bottleneck.

\paragraph{Where the MTP cycle idles} The idle time has one main cause~\ev{profile-host}. Outside graph replays the GPU waits about 3~ms per profiled cycle at every batch size, and the draft and verify graphs are launched only 44--46~\textmu s and 380--411~\textmu s before the GPU reaches them. NVTX ranges on the speculative worker's host functions and samples of the scheduler put 1.2--1.5~ms of exposed host time per cycle in the attention planning for the verify pass (FlashInfer's prefill planner, reached through the verify preparation and the graph runner's batch loading, with blocking device-to-host copies), of which only 14\% overlaps GPU work at $B=1$, and another 0.5~ms in the draft pass's planning; the host work of the draft-extend pass is 97\% hidden behind the verify graph. Subtracting the traced GPU-busy time from the cycle measured without a profiler leaves about 21\%, 20\% and 15\% of the cycle idle at $B=1$, 8 and 32 (Appendix~\paperref{app:profile} of the paper).

\begin{table}[htbp]
\small\centering
\begin{tabularx}{\linewidth}{@{}YP{3.1cm}rr@{}}\toprule
Cost & Configuration & $f$ & Ceiling\\\midrule
FP32 recurrent state read and written every step & plain, $B=128$ ($B=32$) & 41.6\% (18.0\%) & $1.71\times$ ($1.22\times$)\\
GPU idle in the cycle, 35.5 of its 37.3 points in host-bound gaps outside the graph replays (attention planning) & MTP, $B=1$ & 37.3\% ($\approx$21\%) & $1.59\times$ ($\approx1.27\times$)\\
GDN verify kernel and state commit, of which only the per-position FP32 state writes and the commit copy are removable & MTP, $B=32$ & 27.2\% & $1.37\times$ (whole)\\
Weight GEMMs below the head GEMM's bandwidth (excess time only) & plain, $B=32$ ($B=1$) & 15.5\% (13.2\%) & $1.18\times$ ($1.15\times$)\\
Three full-vocabulary draft heads per cycle & MTP, $B=1$ ($B=32$) & 12.9\% (8.5\%) & $1.15\times$ ($1.09\times$)\\
The target head chain & plain, $B=1$ ($B=128$) & 10.3\% (6.1\%) & $1.11\times$ ($1.07\times$)\\
65 RMSNorm launches per step & plain, $B=8$ & 4.8\% & $1.05\times$\\
Gaps between graph nodes & plain, $B=1$ & 3.0\% & $1.03\times$\\
Eager sampling and bookkeeping after each replay & plain, $B=1$ & 2.6\% & $1.03\times$\\\bottomrule
\end{tabularx}
\caption{Measured costs and their Amdahl ceilings $1/(1-f)$, with $f$ the share of the step or cycle in the stated configuration. For the idle time the parenthesized values are the share derived without the profiler and its ceiling. Measured~\ev{profile-attr,profile-host}.}
\label{tab:profile-ranked}
\end{table}

\subsection{Ceilings under a stated execution model}\label{app:ceilings}\label{sec:stack-ceilings}
A decode step streams the linear weights once whatever the batch size, $W=8.41$~GB, and everything else is per request: 100.7~MB of FP32 recurrent state read and written, and 32~KiB of attention KV per context token (Table~\paperref{tab:bytes} of the paper). With $F$ the floating-point operations per token (two per weight), $s$ the per-request bytes per step, $\mathrm{BW}$ the measured HBM read bandwidth and $P$ an achievable tensor-core rate, and with each layer's weight GEMMs and its per-request kernels running in turn, as they do at the pin, the step time at batch $B$ is bounded below by
\begin{equation}
t(B)\ge\max\!\left(\frac{W}{\mathrm{BW}},\,\frac{BF}{P}\right)+\frac{Bs}{\mathrm{BW}},\label{eq:stack-floor}
\end{equation}
since in that order the weight stream and the GEMM arithmetic can overlap each other while the per-request traffic adds to them. An engine that overlapped the memory-bound per-request kernels with the batch GEMMs would be bounded only by
\begin{equation}
t(B)\ge\max\!\left(\frac{W+Bs}{\mathrm{BW}},\,\frac{BF}{P}\right),\label{eq:stack-floor-overlap}
\end{equation}
whose large-batch limit is $\min(\mathrm{BW}/s,\,P/F)$. Equation~\eqref{eq:stack-floor} models the engine as it runs; Equation~\eqref{eq:stack-floor-overlap} is the overlapped roofline at the same inputs, for an engine that splits the batch to overlap the two kinds of kernel. Neither is a bound for every engine on this machine. The bandwidth input is the 3.79~TB/s measured over the head's footprint~\ev{profile-bw}, and $P$ is assumed to be 70\% of the BF16 and 60\% of the FP8 datasheet peaks, assumptions stated so that the ceilings can be recomputed. Two consequences follow. Under Equation~\eqref{eq:stack-floor}, as $B$ grows the throughput ceiling approaches $1/(F/P+s/\mathrm{BW})$, set by the per-request state and the arithmetic rather than by the weights, so halving the state's bytes or writing it less often raises it while a faster head does not; and even with every byte removed, the arithmetic term caps an FP8 engine, so no single lever gives an order of magnitude at high concurrency. At batch one the weights bound the step, and the draft head and the target's weight bytes are the terms to cut.

Table~\ref{tab:stack-ceilings} and Figure~\ref{fig:stack-ceilings} evaluate the floor for stacks of levers~\ev{ceilings}; they are derived, not measured. At batch one the floor is 2.25~ms per token (445 tokens/s), of which the weights alone account for 2.22~ms; plain decoding measured 280 tokens/s at concurrency one in a single served run~\ev{levers}, 63\% of that rate, and 289 tokens/s in the profile's fixed-batch windows~\ev{profile-attr}. At large batch the plain engine's ceiling under Equation~\eqref{eq:stack-floor} is 23.7k tokens/s, because each request moves 114~MB per step and needs 8.4~GFLOP. Halving the state bytes, writing the state less often and moving weights and cache to FP8 raise the ceiling to between 34k and 79k tokens/s. The overlapped roofline is higher, 32.1k tokens/s for the plain engine and 141.2k for the full lossy stack at $B=2{,}048$ (the plain engine's limit as $B$ grows is 33.3k), but under neither model does any stack reach ten times the plain engine's ceiling: 8.4~GFLOP per token caps an FP8 engine near 140k tokens/s at the assumed rate even with every byte removed (about 235k at the datasheet FP8 peak, of which the assumed rate is 60\%; derived from the same constants).

\begin{table}[htbp]
\small\centering
\begin{tabular}{@{}lrrr@{}}\toprule
& MB per request & \multicolumn{2}{c}{Tokens/s at $B=2{,}048$}\\\cmidrule(l){3-4}
Lever stack & per step & layer-serial & overlapped\\\midrule
Plain, FP32 state & 114.0 & 23.7k & 32.1k\\
16-bit state & 63.6 & 34.6k & 55.9k\\
Replayed state writes, FP32 & 66.8 & 33.6k & 53.5k\\
Replayed state writes, 16-bit state & 40.0 & 44.0k & 82.3k\\
Replayed state writes, int8 state & 26.7 & 52.1k & 82.3k\\
Replayed, 16-bit state, FP8 weights and KV & 34.6 & 61.7k & 102.6k\\
Replayed, int8 state, FP8 weights, KV and head & 21.2 & 78.9k & 141.2k\\
MTP (accepted length 3), stock verify, FP32 & -- & 19.5k & --\\
MTP (accepted length 3), replayed verify, FP32 & -- & 34.3k & --\\\bottomrule
\end{tabular}
\caption{Derived large-batch throughput ceilings under the layer-serial model of Equation~\eqref{eq:stack-floor}, which assumes that per-request kernels do not overlap the weight GEMMs, and under the overlapped roofline of Equation~\eqref{eq:stack-floor-overlap}, at a 334-token context, with the measured read bandwidth of 3.79~TB/s and the tensor-core rate assumed at 70\% (BF16) or 60\% (FP8) of the datasheet peak. The MTP rows are for the layer-serial model only. Not measurements~\ev{ceilings}.}
\label{tab:stack-ceilings}
\end{table}

\begin{figure}[htbp]
\centering
% Derived throughput ceilings per lever stack and single served runs (research notes,
% fig:stack-ceilings); data: evidence/moonshot/ceilings.csv and lever_sweeps_quick.csv.
% Every series is plain decoding (no drafter), so the layer-serial model is ink and the overlapped
% roofline, the alternative model, grey (style.tex), with one dash pattern per lever stack. The
% ceilings are derived (lines without markers). The served runs are measured (filled markers,
% single runs), offset by a factor 1.12 left and right of their concurrency so that both stay
% visible.
% Direct labels, no legend: each ceiling is named right of the plot with a leader from its end at
% B = 2048, the labels fanned 0.3 cm apart in the order of the lines' ends (23.7k, 32.1k, 34.6k,
% 44.0k, 78.9k, 141k tokens/s, checked against ceilings.csv); the served runs are named at c = 1,
% FP16 state right of its square and plain under it, with a leader to its circle.
\datafigure{../evidence/moonshot/ceilings.csv}{%
\def\ceildata{../evidence/moonshot/ceilings.csv}%
\def\served{../evidence/moonshot/lever_sweeps_quick.csv}%
\vpreadcsv{\ceiltable}{\ceildata}%
\vplookuptwo{\ceiltable}{stack}{"baseline (FP32 state)"}{batch}{2048}{tokens_per_s_ceiling}{\ceilA}%
\vplookuptwo{\ceiltable}{stack}{"16-bit state (FP16 or BF16)"}{batch}{2048}{tokens_per_s_ceiling}{\ceilB}%
\vplookuptwo{\ceiltable}{stack}{"ReplaySSM + 16-bit state"}{batch}{2048}{tokens_per_s_ceiling}{\ceilC}%
\vplookuptwo{\ceiltable}{stack}{"ReplaySSM + int8 state + FP8 W8A8 + FP8 KV + FP8 head"}{batch}{2048}{tokens_per_s_ceiling}{\ceilD}%
\vplookuptwo{\ceiltable}{stack}{"baseline (FP32 state)"}{batch}{2048}{overlapped_tokens_per_s_ceiling}{\ceilE}%
\vplookuptwo{\ceiltable}{stack}{"ReplaySSM + int8 state + FP8 W8A8 + FP8 KV + FP8 head"}{batch}{2048}{overlapped_tokens_per_s_ceiling}{\ceilF}%
\begin{tikzpicture}
\begin{axis}[width=.62\linewidth,height=6.0cm,xmode=log,log basis x=2,ymode=log,xmin=0.8,xmax=2400,ymin=100,ymax=2e5,
  xminorticks=false,yminorticks=false,clip=false,
  xtick={1,8,32,128,512,2048},xticklabels={1,8,32,128,512,2048},
  ytick={100,1000,10000,100000},yticklabels={100,1k,10k,100k},
  xlabel={Batch size $B$ (lines) or client concurrency $c$ (points)},ylabel={Output tokens/s},grid=major]
\addplot[thick,cstock,derived,discard if not={stack}{"baseline (FP32 state)"}] table[x=batch,y=tokens_per_s_ceiling,col sep=comma]{\ceildata};
\addplot[thick,cstock,dashed,derived,discard if not={stack}{"16-bit state (FP16 or BF16)"}] table[x=batch,y=tokens_per_s_ceiling,col sep=comma]{\ceildata};
\addplot[thick,cstock,dashdotted,derived,discard if not={stack}{"ReplaySSM + 16-bit state"}] table[x=batch,y=tokens_per_s_ceiling,col sep=comma]{\ceildata};
\addplot[thick,cstock,densely dotted,derived,discard if not={stack}{"ReplaySSM + int8 state + FP8 W8A8 + FP8 KV + FP8 head"}] table[x=batch,y=tokens_per_s_ceiling,col sep=comma]{\ceildata};
\addplot[thick,ccontext,derived,discard if not={stack}{"baseline (FP32 state)"}] table[x=batch,y=overlapped_tokens_per_s_ceiling,col sep=comma]{\ceildata};
\addplot[thick,ccontext,densely dotted,derived,discard if not={stack}{"ReplaySSM + int8 state + FP8 W8A8 + FP8 KV + FP8 head"}] table[x=batch,y=overlapped_tokens_per_s_ceiling,col sep=comma]{\ceildata};
\addplot[only marks,mark circle,cstock,discard if not={config}{plain}] table[x expr=\thisrow{concurrency}/1.12,y=y_tok_s_gpu,col sep=comma]{\served};
\addplot[only marks,mark square,cstock,discard if not={config}{plain+fp16_state}] table[x expr=\thisrow{concurrency}*1.12,y=y_tok_s_gpu,col sep=comma]{\served};
% ceiling labels right of the plot, fanned 0.3 cm apart, leaders from the lines' ends at B = 2048
\draw[leader] (axis cs:2100,\ceilF) -- (axis cs:3000,141000) (axis cs:2100,\ceilD) -- (axis cs:3000,82000)
  (axis cs:2100,\ceilC) -- (axis cs:3000,48000) (axis cs:2100,\ceilB) -- (axis cs:3000,28000)
  (axis cs:2100,\ceilE) -- (axis cs:3000,16400) (axis cs:2100,\ceilA) -- (axis cs:3000,9600);
\node[annot,anchor=west] at (axis cs:3050,141000) {replayed, int8 state, FP8, overlapped};
\node[annot,anchor=west] at (axis cs:3050,82000) {replayed, int8 state, FP8, layer-serial};
\node[annot,anchor=west] at (axis cs:3050,48000) {replayed, 16-bit state, layer-serial};
\node[annot,anchor=west] at (axis cs:3050,28000) {16-bit state, layer-serial};
\node[annot,anchor=west] at (axis cs:3050,16400) {plain, FP32 state, overlapped};
\node[annot,anchor=west] at (axis cs:3050,9600) {plain, FP32 state, layer-serial};
% served runs at c = 1 (circle at 1/1.12, square at 1.12, both near 280 tokens/s)
\node[annot,anchor=west] at (axis cs:1.3,284) {FP16 state, served (single run)};
\draw[leader] (axis cs:0.95,255) -- (axis cs:1.25,185);
\node[annot,anchor=west] at (axis cs:1.3,180) {plain, served (single run)};
\end{axis}
\end{tikzpicture}}{Derived throughput ceilings per lever stack against batch size, with measured served throughput.}%

\caption{Derived throughput ceilings per lever stack against batch size under the layer-serial model of Equation~\eqref{eq:stack-floor} and, for two stacks, the overlapped roofline of Equation~\eqref{eq:stack-floor-overlap}, at the measured bandwidth and assumed tensor-core rate (lines~\ev{ceilings}), and served output throughput of single benchmark runs at client concurrency 1, 32 and 128 (points~\ev{levers}; the two series are offset slightly left and right of their concurrency so that both stay visible). The gap between a line and a point at the same concurrency is overhead above the layer-serial floor; at high concurrency it includes the server's streaming path (Appendix~\paperref{app:serving-hazards} of the paper).}
\label{fig:stack-ceilings}
\end{figure}

Repeated engine-only and served measurements up to $B=1{,}024$ for each state and weight format are \pend{moonshot-repeat}.

\subsection{Speculation and the recurrent state}\label{app:spec-state}\label{sec:stack-state}
The hybrid architecture changes the economics of speculation. A verify pass over $D$ draft tokens must be able to roll back to any accepted prefix, and SGLang's Triton verify kernel does so by writing one full FP32 state per draft position into an intermediate buffer and copying the accepted one back at commit. Per request and cycle that is one state read, $D$ state writes and a read and write at commit: 352~MB for $D=4$, three and a half plain steps (derived from the code path and checked against the trace). The two kinds of traffic scale differently with the batch. Speculation saves weight bytes: one verify pass streams the weights once for several committed tokens, and that stream is shared by every request in the batch, so its cost per output token falls as $1/B$. The snapshot writes are per request and per draft position, so their cost per output token does not fall with the batch at all. At batch one the weight saving dominates; at high batch, where the state dominates, speculation's bandwidth advantage shrinks and can turn into a loss. This cost belongs to a recurrent hybrid; an attention-only model rolls back by discarding cache entries and writes no snapshots. Rollback of recurrent state is an active problem: the Mamba-in-Llama work advances one cached state lazily and recomputes after a rejection~\cite{mambainllama}; STree, TreeWY, Bole, GDN Tree-Scan and SpecLA verify trees or chains on SSM and GDN hybrids~\cite{stree,treewy,bole,gdntreescan,specla}, and TreeWY states that its outputs are not bit-identical to the baseline; ReplaySSM caches recent inputs instead of writing the state at every step, equivalent to ordinary decoding up to floating-point error~\cite{replayssm}.

The measurements agree with the mechanism. At $B=32$ the verify kernel runs for 2.54~ms and the commit for 1.23~ms per cycle, 27\% of the cycle, at an implied 3.2 and 2.6~TB/s~\ev{profile-attr}. The byte model, at a context of 700 tokens and an accepted length of 2.8, has MTP moving 44\% fewer bytes than plain decoding per output token at $B=1$, 26\% fewer at $B=32$, 6\% fewer at $B=128$ and 2\% more at $B=256$~\ev{profile-bytes}; the crossover moves with the context, because attention KV is read once per verify pass for several committed tokens, so longer contexts favour speculation. In the profile's fixed-batch windows MTP's throughput advantage over plain decoding fell from $1.33\times$ at $B=1$ to $1.13\times$ at $B=32$ (Appendix~\paperref{app:profile} of the paper). Served single runs on the tuning split show the same turn: MTP with three steps and snapshot verification led plain decoding by $1.64\times$ at concurrency 1 and trailed it from concurrency 32, and replayed verification, which removes the snapshots, raised its throughput at concurrency 128 by 24\% (Table~\paperref{tab:tuning} of the paper). At that point the snapshots alone would be written at about 0.59~TB/s if each draft token's 50.3~MB snapshot is written once per cycle, against 1.35~TB/s of state traffic for plain decoding at the same concurrency under the same assumption (derived from the measured throughput and accepted length)~\ev{tuning}. On the confirmation split, over three sessions with matched plain baselines, every speculative family leads plain decoding through concurrency 32 and trails it from 48, with accepted length flat across concurrency (3.24 to 3.28 tokens per cycle for MTP, 4.66 to 4.81 for DFlash at block 8), so the loss is not a fall in acceptance. At concurrency 128 MTP with replayed verification delivers 0.86 of plain decoding's throughput, and with SGLang's snapshot verification 0.68 (one session)~\ev{bench-frontier}.

Bytes do not explain the whole cycle at high concurrency~\ev{spec-cycle}. At concurrency 128 the tuning points imply an MTP cycle of 43.5~ms with snapshot verification and 35.0~ms with replayed verification (128 requests times 3.26 tokens per cycle, over the measured throughput), against 9.5~ms per plain step. The implied cycle is derived from client throughput, so it includes the prefill passes of arriving requests and the benchmark's ramp and drain as well as the decode cycle. Derived from bytes and FLOPs at 3.8~TB/s, 650~TFLOP/s and a context of about 400 tokens, its parts add up to about 21.7 and 13.2~ms: the verify GEMMs over 512 tokens take 6.6~ms, the three draft steps about 1.7~ms, verify attention and small kernels about 1.5~ms, and the GDN state traffic 11.9~ms with snapshots (352~MB per request) or 3.4~ms with replay (about 100~MB). Replay saved about 8.5~ms per cycle, which the 252~MB per request of snapshot and commit traffic it removes accounts for, but the implied cycle is 2.0 to 2.7 times the sum of its derived parts: most of it is time the byte model does not see, in the decode cycle itself or in the prefill, ramp and drain, which this arithmetic cannot separate.

SGLang's replay path for speculation stores each draft token's inputs and folds the accepted prefix into the state at commit, reading and writing the state once per cycle, which would make speculation cheaper in state bytes per token than plain decoding; trees do not get this benefit yet, because the replay path handles only chains and a tree keeps one full state per node~\cite{treewy,gdntreescan}. The block drafter is more exposed than the MTP layer: a 16-token DFlash block writes about 805~MB of states per request per cycle (Section~\ref{app:drafter-model}).

Verifying a draft block in parallel over its positions (P7b), instead of through the sequential recurrence, would make wide blocks cheap only with a fast path that is both faster and close enough to the recurrent outputs to be certified. SGLang's chunked GDN kernel is neither at one request~\ev{fast-verify}. On one layer, with synthetic activations and the checkpoint's gates, its outputs over a block of 2 to 16 tokens differ from the packed decode's in 52--55\% of BF16 output words, and it takes 815--990~\textmu s per layer at block widths 4 to 256, against 75--378~\textmu s for the recurrent verify kernel with its per-position snapshots. A purpose-built block-parallel kernel is untested.

State precision is the cheapest lever and the most dangerous. DAMP reports reasoning accuracy of 85.5 with FP32 recurrent state, 84.6 with FP16, 79.7 with BF16 and below 30 with FP8 or int8 on another GDN model~\cite{damp}, and LeapQuant recovers near-FP32 accuracy with 8-bit state by re-quantizing once per 16-token window~\cite{leapquant}. SGLang accepts FP16 state (\code{-{}-mamba-ssm-dtype float16}). MTP with and without replayed verification at concurrency 128 to 1,024, and the logit-probe KL, divergence rate and GSM8K accuracy for FP32, FP16, BF16 and FP8 state with and without replay, are \pend{moonshot-state}.

\subsection{How much recurrent state an exact computation needs}\label{app:minimal-state}\label{sec:minimal-state}\label{sec:stack-state-need}
The dense FP32 state is sufficient for exact decoding; nothing yet shows that it is necessary. Two questions should be kept apart: whether less information could be retained, and whether the state has to be materialized as often as it is.

For the first, fix one GDN head of one layer and treat the layer's inputs as given. Along any permitted history the state evolves as $S_t=S_{t-1}A_t+B_t$ with $A_t=\alpha_t(I-\beta_tk_tk_t^\top)$ and $B_t=\beta_tv_tk_t^\top$, which depend on the inputs and not on the state, and the head reads $o_t=S_tq_t$. A state difference $\Delta S$ at time $t$ therefore propagates linearly, to $\Delta S\,A_{t+1}\cdots A_{t+j}$, and is observed as $\Delta S\,A_{t+1}\cdots A_{t+j}q_{t+j}$.
\begin{proposition}[What an exact state must retain]\label{prop:minimal-state}
Let $\mathcal S$ be the set of states reachable over permitted histories, $\mathcal D=\operatorname{span}\{S-S':S,S'\in\mathcal S\}$, and $\mathcal N$ the set of differences $\Delta S$ with $\Delta S\,A_{t+1}\cdots A_{t+j}q_{t+j}=0$ for every permitted continuation and every $j\ge0$. Two reachable states produce the same outputs under every permitted continuation if and only if their difference lies in $\mathcal N$. If $\mathcal S$ contains a relatively open subset of its affine hull, then a linear map that separates every pair of reachable states with different outputs has rank at least $\dim\mathcal D-\dim(\mathcal D\cap\mathcal N)$ on $\mathcal D$.
\end{proposition}
\begin{proof}
Because $B_t$ does not depend on the state, the output difference under a continuation is the linear expression above, and it vanishes for every continuation exactly when $\Delta S\in\mathcal N$. For the bound, take $d\in\mathcal D$ with $\phi(d)=0$. By the hypothesis there are a reachable $S$ and an $\varepsilon>0$ with $S+\varepsilon d$ reachable; $\phi$ gives them the same image, so they must have the same outputs, and $\varepsilon d\in\mathcal N$. Hence $\ker\phi\cap\mathcal D\subseteq\mathcal N\cap\mathcal D$, and the rank of $\phi$ on $\mathcal D$ is at least $\dim\mathcal D-\dim(\mathcal D\cap\mathcal N)$.
\end{proof}
The hypothesis matters: if only one state is reachable, $\dim\mathcal D=0$ and a representation with no dimensions is exact. Treating every input sequence as permitted is conservative, since in the full model a layer's outputs influence later tokens; it is the notion that makes a reduction safe for any request. Exact witnesses show that neither $\mathcal D$ nor $\mathcal N$ is favourable without structure in the keys~\ev{state-structure}. Writes with unit coordinate keys and gates inside $(0,1)$ reach any matrix, so the reachable set can fill the whole state space and $\mathcal D$ with it. A direction the current query cannot see is not unobservable: with $\alpha=\beta=\tfrac12$, $k=(\tfrac35,\tfrac45)$, $q=e_1$ and $\Delta S=ue_2^\top$ for any value vector $u$, the observed difference $\Delta S\,q$ is zero now, but one step later it is $\Delta S\,Aq=-\tfrac{3}{25}u$. Each transition is a contraction, not an erasure: $\det A_t=\alpha_t^d(1-\beta_t\norm{k_t}^2)>0$ when $\beta_t\norm{k_t}^2<1$, so no step forgets a direction exactly. Transitions with non-orthogonal keys do not commute, so there is no common basis in which they are diagonal. And because a gated RMSNorm precedes the output projection, the null space of the projection is not unobservable: two outputs that $W_o=(1,0)$ cannot tell apart, $(1,1)$ and $(1,7)$, project to $1$ and $\tfrac15$ after normalization.

One structural reduction is exact. If every key lies in the column space of a fixed orthonormal $U\in\R^{d_k\times r}$ and $S_0=Z_0U^\top$, then $S_t=Z_tU^\top$ with
\[Z_t=\alpha_tZ_{t-1}+\beta_t\bigl(v_t-\alpha_tZ_{t-1}\kappa_t\bigr)\kappa_t^\top,\qquad\kappa_t=U^\top k_t,\]
the same recurrence in dimension $r$, since $U^\top A_t(I-UU^\top)=0$ and the write $\beta_tv_tk_t^\top$ stays in the form exactly when $k_t\in\operatorname{col}(U)$. The keys are produced by a short convolution, a SiLU and a normalization, and SiLU raises rank ($\operatorname{SiLU}(\log2\,[[1,2],[2,4]])$ has rank two), so a low-dimensional input does not imply low-dimensional keys; whether real keys are confined to a subspace is a measurement, to be certified by exact or modular rank rather than an SVD threshold \pend{moonshot-rank}. With $U$ fixed per layer and head and shared across requests, the per-request state is $d_vr$ values instead of $d_vd_k$, a saving for every $r<d_k$; a per-request low-rank factorization $S=AB^\top$ instead stores $d_vr+rd_k$ values, which beats the dense state only below $r=64$ at $d_v=d_k=128$ and halves it only at $r\le32$ (derived). Floating point adds a further obstruction: an exact rational reduction does not preserve bitwise results. The split form of the recurrence changed 1,481 of 1,600 FP32 outputs (Appendix~\paperref{app:contracts} of the paper), deferring a decay reassociates FP32 roundings (one ulp apart for $\alpha_1=\alpha_2=0.9$, $x=1.3$), and rounding an exactly rank-one BF16 outer product to BF16 gives a $32\times32$ matrix of exact rank 28~\ev{state-structure}.

Errors that do enter the state stay bounded as long as each transition is non-expansive: $(I-\beta kk^\top)$ has spectral norm at most 1 exactly when $0\le\beta\norm k^2\le2$ (its eigenvalues are 1 and $1-\beta\norm k^2$), so a perturbation $\eta_t$ injected at each step obeys $\norm{E_t}\le|\alpha_t|\norm{E_{t-1}}+\norm{\eta_t}$; with normalized keys and $\beta\in(0,1)$ the condition always holds~\ev{contracts}.

The second question has a better answer. A decoder that keeps a dense anchor and the recent actual FP32 operands of each step, and replays them in the reference kernel's order, reproduces every state word bit for bit by induction, while reading the dense state once per step and writing it once every $m$ steps. That reduces information movement, not information, and it is what replay-based rollback already does~\cite{replayssm,gdncode}. The paragraphs below treat it as a lever with a pre-registered test (P4), and the opposite idea, sharing exact computation across requests (P5).

\paragraph{Writing the recurrent state less often (P4, hypothesis)} Rounding-preserving live replay is exact and cuts dense-state traffic per step from $2D$ (read and write) to $D+D/m$ when the state is flushed every $m$ steps, 0.625 of the original at $m=4$, for a few percent more storage (derived). If the recurrent kernel is a fraction $f$ of a decode step, a kernel twice as fast gives at most $1/(1-f/2)$, removing it entirely $1/(1-f)$, and the $m=4$ traffic cut alone $1/(1-0.375f)$. The measured share at $B=128$ and a 418-token mean context is $f=0.416$ (Table~\paperref{tab:profile-plain} of the paper), which gives $1.26\times$, $1.71\times$ and $1.18\times$ (derived from a measured share). The pre-registered test has two parts. A kernel check compares every state word and GDN output of the replayed and dense paths, one layer at a time, and one differing word rejects bit-exactness (below). The served timing decodes 512 tokens greedily at batch 128 from 2,048-token prompts with FP32 state and no speculation, with the threshold at $1.10\times$; through the server only the output probe below checks exactness. The served timing's design, metric and decision rule were declared before its first run, and every later change before the run it governs~\ev{p4-declared}. Dense decoding and replay at $m=4$ run as four server pairs in the order A B B A A B B A with identical memory pools. The primary metric is the server's decode rate at exactly 128 running requests: the harmonic mean of the scheduler's generation-throughput windows (one per 40 decode passes) that show 128 running requests on that line and the previous one, with no prefill batch between them, inside the AIPerf profiling phase, which excludes warm-up, ramp and drain; fewer than 8 such windows in any arm voids the run. It measures decode only, which is what the lever changes, and client throughput, which includes the 2,048-token prefills, is reported beside it. The four log ratios give a 95\% $t$ interval, and the claim is rejected if its upper end is below $1.10$, supported if its lower end reaches $1.10$, and inconclusive otherwise; a supported result is a served decode throughput at 128 running requests, not an end-to-end speedup. A server output probe compares greedy tokens and top-20 log-probabilities at concurrency 1 with dense decoding, with a second dense server as the noise control: a difference refutes end-to-end exactness unless the control also differs from the first dense server, in which case the probe is undecided; a pass does not establish exactness; and a refuted probe leaves P4's exact claim unsupported whatever the throughput interval. A run that fails its completeness, workload, pool, environment or provenance checks is void and its numbers are not reported. One check was corrected after the first run and before any verdict was computed: the sent prompts are now matched by the hashes of their texts, because the benchmark's request ids differ from the declared ones for this workload, a check no weaker than the declared one. Three attempts have been void, because in every arm at most 127 of the 128 requests were admitted. The amendments declared after the first two, a pinned attention cache of 655,360 tokens, 132 pinned recurrent-state slots and an admission preflight that must show 128 running requests, did not change that: in the third the attention cache was at 40\% use and 127 of 132 recurrent-state slots were taken. Read from the engine source, the admission check counts a request whose last prefill chunk runs in the pass twice: the request-table row it keeps from its first chunk lowers the admission limit, and it also counts against that limit as one of the pass's requests. A wave therefore stops one short of the running limit whenever its last pass carries a chunk tail, and the 128th request runs only after the wave finishes. Every one of the 56 admission waves in the three attempts' server logs had such a tail and stopped at 127, as predicted~\ev{p4-declared}. Those waves test the prediction only at that one limit and one chunk; the discriminating test is the amended admission preflight. Before any further run the running limit was raised to 129, with the client still at 128 concurrent requests, so that at most 128 requests are ever in the server and the prediction is a peak of 128 running in both arms. That preflight is queued. No served result exists. SGLang's own replay path qualifies only if it passes the bit-exactness check. The replay kernel is committed as an engine patch with a one-layer test that compares every output and state word against the stock decode kernel (\code{tests/test\_gdn\_exact\_replay.py}~\ev{moonshot-patches}).

\paragraph{P4 at kernel level} The replay is bit-exact at kernel level~\ev{p4-replay}. Run side by side with SGLang's packed decode for 48 steps at batch 8, on synthetic activations with the checkpoint's gates, with rows flushing at staggered phases and the reconstructed state written out and compared at every step, it differs in none of 201,326,592 state words and none of 1,572,864 output words, at $m=4$ and 16 on the first GDN layer and at $m=4$ on layer 20. SGLang's own replay path differs on the same inputs in 465,102 and 543,225 of the 1,572,864 output words (30\% and 35\%, layers 0 and 20), so it does not qualify. Timed on one layer under CUDA-graph replay with a cold L2, the kernel is $1.215\times$ faster than the packed decode at batch 128 and $1.227\times$ at batch 256 with $m=4$, not the $1.6\times$ the bytes allow: at batch 128 it takes 89, 112 and 129~\textmu s on steps that replay zero, one and two earlier updates and 201~\textmu s on the step that replays three and writes the anchor, against 161.5~\textmu s for the packed decode; with $m=8$ and 16 it is no faster than the packed decode. With the recurrent kernel at 41.6\% of a batch-128 step, a kernel $1.22\times$ faster predicts about $1.08\times$ for the decode step, and less with the pre-registered 2,048-token prompts, where attention takes a larger share: below the $1.10\times$ the test requires (derived). The check covers one layer at a time on synthetic activations; the server output probe and the served paired timing are \pend{moonshot-p4}.

\paragraph{Sharing computation across requests (P5, rejected)} Factoring hidden or recurrent state across unrelated requests does not survive analysis. An exact factorization with a dense residual costs the full projection, so it saves only if $r(d+B)+z<Bd$ with a zero or sparse residual; SiLU and gating raise rank; full-rank pairwise state differences force private factors at least as large as the dense state; one gated-delta step takes four identical states to batch rank four~\ev{state-structure}; online factorization costs of order $B^2$ times the update; and reassociation breaks equality with the reference. The one exact candidate is structural: the first GDN layer's input projections (12,352 outputs) depend only on the current token, so they could be read from a precomputed per-token table (5.71~GiB for the whole vocabulary) with convolution and recurrence kept private. Its derived ceiling is about 0.75\% of the dense projection work per token, and at batch 128 even a 99\% table hit rate would leave about 72\% of batches with a miss. The pre-registered criterion rejects the table unless the stage's measured share of a decode step allows a 1\% end-to-end gain. It does not: layer 0's input projections take 0.60\%, 0.57\%, 0.51\% and 0.33\% of a plain step at $B=1$, 8, 32 and 128, a ceiling of $1.006\times$~\ev{profile-p5}, so P5 is rejected by its criterion as well as in analysis.

\subsection{Levers by axis}\label{app:levers}\label{sec:stack-levers}
Table~\ref{tab:axes} lists the levers, each built as a separate server flag or environment variable so that any subset composes. A lever is \emph{exact} when it leaves the target's arithmetic unchanged, so every emitted token is the target's decision at its own batch shape under the stock-kernel contract, with its measured noise floor; drafting levers are exact because the target verifies every token. A lever that changes the target's arithmetic is classed by its greedy outputs against a matched stock reference (Appendix~\paperref{app:serving-protocol} of the paper): \emph{exact up to rounding} if every first divergence is at an exact tie, within one BF16 step or a near tie within two steps, the classes of the stock engine's own batch-shape divergences (Appendix~\paperref{app:divergence} of the paper), and \emph{lossy} otherwise, in which case it carries a measured quality cost. SGLang's replay path for plain decoding is exact up to rounding by that comparison~\ev{bench-equality}, although it is not bit-exact: its decode kernel reconstructs the state on tensor cores from BF16-cast operands and persists that rounding at every flush (code reading at the pin), and on synthetic activations its outputs differ from the packed decode's in 30 to 35\% of words (Section~\ref{app:minimal-state}). The bit-exact alternative is the rounding-preserving replay (P4).

The first single-run measurements on the benchmark harness (plain arm with the prefix cache on, capacity 128) separate the levers by where they act~\ev{levers}. FP16 recurrent state raised served throughput at concurrency 128 from 13.5k to 16.7k tokens/s ($1.24\times$) and changed nothing at concurrency one, which is what a per-request byte term predicts; its quality cost is \pend{moonshot-state}. These runs streamed one token per chunk, and at that concurrency the streaming path rather than the GPU can limit a plain arm (Appendix~\paperref{app:serving-hazards} of the paper), so the ratio may understate the lever; a paired rerun is \pend{moonshot-repeat}. FP8 KV did nothing at these short contexts ($0.99\times$). FP8 W8A8 weights could not use their CUTLASS GEMM: the one in the installed aarch64 \code{sgl-kernel} aborts at run time because that binary was not compiled for \code{sm\_90a}, as its logged error states (Hopper's asynchronous warp-group MMA instructions exist only for that target~\cite{ptxisa,hopper}), and the Triton route showed no consistent gain ($0.96$--$1.06\times$ across concurrencies), a negative result until a different GEMM route is wired. A first engine-only pass of decode steps at batch 512 took 26.9~ms with FP32 state, 23.0~ms with BF16 and 22.5~ms with FP16~\ev{decode-steps}; it is noisy at small batch sizes and is to be repeated. All of these are single runs without repeats.

Two levers need engine changes. Draft-vocabulary methods restrict the draft head to a subset of rows~\cite{frspec,vocabtrim,specvocab,dynaspec,microspec}, and FastMTP compresses only the draft vocabulary while keeping full-vocabulary verification~\cite{fastmtp}; they change only the drafter, so the target's decision is unchanged under target-only or point-mass verification. SGLang's hot-vocabulary draft head (\path{--speculative-token-map}, after FR-Spec) does the wrong thing on Qwen3.5 as read at the pin: the MTP layer's head is the tied embedding, and \code{set\_embed\_and\_head} ignores the truncated head in that case, so the drafter scores all 248,320 rows while the worker remaps indices as if it had scored the hot subset; the patch unties a reduced head for the drafter only. A hot map ranked by token frequency on the tuning split's greedy outputs covers 97.5\% of the 171,968 output tokens of a plain sweep on the confirmation split with 22,936 rows (every token seen in the tuning outputs), 96.3\% with 16,384 and 89.1\% with 4,096~\ev{token-map}. The second change adds a relaxed greedy rule for chains, accepting a draft token whose target logit is within $g$ of the maximum; SGLang's existing relaxed thresholds act only when sampling, and the rule is lossy by construction, so its quality curve is the measurement~\cite{marginsnotwindows}. Both changes are committed as engine patches with tests that run where SGLang and a GPU are available (\code{tests/test\_moonshot\_levers.py}~\ev{moonshot-patches}); their measured effect and the relaxed rule's quality curve are \pend{moonshot-patches}.

The host gap of the speculative cycle at low concurrency (Appendix~\paperref{app:profile} of the paper) has a configuration-level lever~\ev{host-levers}. With Triton attention for the target and the draft, neither attention backend needs host-side sequence lengths, so the overlap scheduler stops copying them to the host and FlashInfer's planning leaves the cycle. MTP with three steps then gains $1.18\times$ in per-request rate at concurrency 1 and $1.14\times$ in throughput at concurrency 4, with unchanged accepted length; Triton attention for the draft alone gains $1.08\times$ and $1.06\times$ (single runs). The draft-only change is exact, because it changes only which tokens are proposed. Triton attention for the target changes the target's arithmetic and is exact up to rounding~\ev{bench-equality}. SGLang's overlapped planning stream fails at the first verify on this hybrid model, whose GDN backend does not implement a call the stream needs.

\paragraph{Removing the blocking reads from attention planning} With FlashInfer attention the host gap has an exact lever as well~\ev{hostgap}. Without a profiler, stock SGLang leaves the GPU idle 1.39 to 1.81~ms of each cycle of the tuned MTP arm at $B=1$ to 128, 20\% of the cycle at $B=1$ and 9\% at $B=128$ (derived: the untraced stock cycle less the GPU busy time per cycle in a host trace of the same configuration; stock and patched traces agree on that busy time within 0.02~ms at $B\le 64$). For a DFlash cycle at block 8 the estimate is 1.36 to 1.46~ms, 21\% at $B=1$ and 9\% at $B=32$ (from the patched engine's trace, assuming the same GPU work). In host traces the largest site is FlashInfer's planning of the verify pass for MTP, which reads a mask size and three index arrays back to the host on every verify, and the drafter's attention planning for DFlash. Engine patches (\path{engine/sglang/patches/hostgap/0001}--\path{0005}) compute those values on the host from lengths the scheduler already holds and leave the plan unchanged. A fourth patch makes that host arithmetic cheaper and uploads the mask offsets from pinned memory, which also removes a last blocking copy in FlashInfer's mask setup. A fifth follows FlashInfer's split-KV rule for NVFP4 caches; it changes nothing with the BF16 cache here, and NVFP4 is untested. On the GPU the plans and a replayed attention graph's output bits matched FlashInfer's own on random batches at every tested size. In the serving engine a validation mode compared at least 8,192 plans of each kind with the stock path and found no difference. Greedy outputs were token-identical to the stock engine in the same configuration, with equal verify-step counts, for all 192 requests (64 prompts at concurrency 1, 8 and 32, top-$k$ 1 chains, prefix cache off, running limit and recurrent-state slots pinned): for MTP with all five patches, and for DFlash with the first three, whose attention cache varied by at most 155 tokens between launches and never bound. Trees, sampling and the prefix cache are not covered. In held decode batches without a profiler the five patches shorten the MTP cycle by 5.9 to 8.8\% (0.50 to 1.20~ms), 36 to 71\% of the idle. The first three alone, which only remove the reads, gave 2.7 to 4.8\% in a session that ran stock first; the five ran patched first. So at low batch most of the idle was host work rather than waiting on a read. Served on the tuned MTP arm, interleaved as stock, patched, patched, stock, the five patches raise per-request rate by 9.4 to 9.6\% at concurrency 1 to 8, 8.9\% at 16 and 7.5\% at 32. All twelve stock--patched pairs fall between $1.067\times$ and $1.101\times$, with unchanged accepted length at concurrency 1 to 16 (one session). What is left is the host path before the draft: the host still needs 1.67 to 1.82~ms (traced) from the previous verify's completion to the next draft's launch, while 0.54 to 0.92~ms of GPU work remains queued. The patches change nothing on the tuned DFlash arm, whose draft uses FA4 attention and whose verify already plans without blocking. With FlashInfer draft attention the DFlash gain (0.7 to 1.7\% in held batches, sequential) is \pend{hostgap-served}.

\paragraph{Backbone GEMMs and RMSNorm} The weight GEMMs' time above the head GEMM's bandwidth and the 65 RMSNorm launches are the fourth and seventh rows of Table~\ref{tab:profile-ranked}. Isolated microbenchmarks under CUDA graphs locate the slow GEMMs and test faster kernels that read the same BF16 weights~\ev{backbone}. Each timed call multiplies every layer's own checkpoint weight by its own input, so the weights stream from HBM as in the model; times are medians of 30 timings of 10 graph replays each, in one exclusive session, include any split-$K$ reduction and the gap to the next kernel, and are compared only within these files. The three projections with 2,560 outputs, the GDN and attention output projections and the MLP down-projection, run as split-$K$ kernels followed by a reduction kernel up to $M=32$ (128 for the down-projection) and reach 49\% (output projections) and 68\% (down-projection) of the measured 3.79~TB/s read bandwidth~\ev{profile-bw} at $M=1$, against 80--87\% for the large projections; disabling cuBLAS's reduced-precision reduction changes neither their time nor any output bit. At $M=1$ SGLang's own Hopper GEMV is the fastest kernel for every projection: it takes 0.77 of cuBLAS's time on the output projections (8.57 against 11.20~\textmu s per call for the GDN one), 0.91 on the down-projection, 0.95--0.97 on the other large projections and 0.51 on the 64-row GDN projection, which cuBLAS runs on a single block. At the pin it is unreachable: \code{-{}-bf16-gemm-backend gemv} loads it, but the unquantized linear layers call the dispatcher only for another backend (code reading; patch \path{engine/sglang/patches/backbone/0001} lets them reach it). At $M=2$--16 a Triton skinny GEMM pays only with programmatic dependent launch, at 0.88--0.92 of cuBLAS's time on the output projections and 0.96--0.99 on the GDN input and MLP projections, with no gain on the attention QKV projection; in this benchmark each call overlaps its successor, which in the model is a different kernel. From $M=64$ cuBLAS, or a cuBLASLt algorithm within 2--5\% of it, is best, except on the down-projection, where PyTorch's cuBLASLt choice takes 0.92 of the time. None of these routes leaves the target's arithmetic unchanged in general: the GEMV and the Triton kernel are bitwise equal to cuBLAS on 99.6 to 100\% of output elements, and the cuBLASLt choice differs from it on the down-projection from $M=16$, while all arms lie within the same distance of a float64 product. If every projection ran at its best isolated time, a plain step's backbone GEMMs would take 197~\textmu s less at $M=1$, at most 5.7\% of the 3,460~\textmu s step, and 92 to 120~\textmu s less at $M=2$--16 (derived from the microbenchmarks, not a served result).

SGLang's packed GDN input projection, one GEMM over the concatenated $12{,}352\times2{,}560$ weight instead of a $12{,}288$-row and a 64-row GEMM, gives every output element bitwise equal to the separate pair's at each tested $M$, the powers of two from 1 to 1,024, with the real weights of all 24 GDN layers and random BF16 inputs, under cuBLAS 13.1 on this GPU~\ev{backbone}. That is a statement about one isolated GEMM, not about served tokens. Against the stock pair, which runs the 64-row GEMM on a side stream, it is as fast up to $M=16$, 3\% slower at 32 and 5 to 13\% faster from 64. The stock fused add-RMSNorm takes 1.47~\textmu s per launch at $M=1$ and 2.0 to 2.2~\textmu s at $M=4$--32 inside a graph of 64 norms. Folding the norm and the SiLU into the next GEMM's prologue reproduces the stock norm bit for bit up to $M=32$ (one element in $10^5$ is one BF16 step away at $M=64$ and 128) and the stock SiLU bit for bit, but as built it makes each layer slower in MLP and GDN layer skeletons, by 12 to 90~\textmu s up to $M=32$, with or without programmatic dependent launch, and by 82 to 543~\textmu s at $M=64$ and 128, far more than the two launches of about 1.5 to 2~\textmu s that it removes. Two further variants, a single-pass norm prologue that scales the product and the SiLU fold alone, also lose: from $M=2$ to 16 they take 103 to 257~\textmu s per layer against 32 to 52~\textmu s for the stock and Triton layers, and at $M=1$ only the SiLU fold ties. In the same skeletons the Triton kernel with programmatic dependent launch saves 1.2 to 2.3~\textmu s per MLP layer and 1.4 to 1.8~\textmu s per GDN layer up to $M=16$, except the GDN layer at $M=4$, and loses at 32. The engine patches that route each projection by a table are committed and off by default. Served, each switch was compared with stock plain decoding on the 320 divergence prompts (256 greedy tokens, top-five log-probabilities, prefix cache on)~\ev{backbone}. With every switch off, and with the packed GDN input projection alone, all 320 prompts return the same output tokens and the same top-five log-probability arrays as stock at concurrency 1; for the packed projection that covers prefill only, since decode at concurrency 1 stays below its 64-row cutoff. SGLang's Hopper GEMV backend diverges at 3.75 per 1,000 compared tokens (164 ties, 8 one-step and 2 near events, none large), against the 3.42 floor at a cap of 16, and the routing table (the Hopper GEMV at $M=1$, the Triton kernel with programmatic dependent launch at $M=2$--16, the packed projection from 64 rows), with pools pinned at a cap of 8, diverges at 3.40 at concurrency 1 and 3.82 at 32, five to six times the 0.63 floor at that cap. Both are exact up to rounding, and no run committed a token that was not its own top-1. In one session of paired serving on tuned plain decoding (order B A A B, the routing table against the same engine with every switch off), the routing table delivered 1.035 and 1.034 times the throughput at concurrency 1, a decode step 116~\textmu s shorter against the 197~\textmu s the microbenchmarks bound; 1.004 in both pairs at concurrency 8, where the decode step is 10~\textmu s shorter against 90 to 98~\textmu s predicted, for a reason not yet traced; 1.003 and 1.001 at 32, about the 0.11\% spread of the two baseline runs, so no claim; and 1.010 in both pairs at 128, where only the packed projection acts and the step is 86~\textmu s shorter against 81 predicted. These are two pairs from one session; serving against tuned MTP and a trace of the GEMM kernels the served engine runs are \pend{backbone-served}.

\begin{table}[htbp]
\small\centering
\begin{tabularx}{\linewidth}{@{}P{2.9cm}YP{1.7cm}P{4.6cm}@{}}\toprule
Axis & Lever & Class & Status\\\midrule
Progress per verify pass & MTP and DFlash drafting, deeper chains, trees where the GDN path allows, n-gram drafting & exact & Chosen in single runs on the tuning split: MTP with three steps, DFlash blocks 16, 8 and 4 by concurrency (Table~\paperref{tab:tuning} of the paper); trees and n-gram drafting not searched; served comparison~\ev{bench-frontier}\\
& Relaxed greedy acceptance within $g$ logits & lossy & \pend{moonshot-patches}\\
Bytes per verify pass & Certified int8 head (Section~\paperref{sec:selfevidence} of the paper) & exact & Head alone (best mode, \modelcons{}): 0.62--0.69 of the stock chain's time up to $M=16$, slower from 128~\ev{kernel-path}; served \pend{bench-cert}\\
& Replayed verification of the recurrent state & exact up to rounding$^\ast$ & MTP, three steps: $0.86\times$ plain decoding at $c=128$ on the confirmation split, against $0.68\times$ with snapshot verification (one session)~\ev{bench-frontier}; outputs rounding-level against stock MTP~\ev{bench-equality}; no detectable GSM8K loss~\ev{bench-quality}; beyond $c=128$ \pend{moonshot-state}\\
& Fold-every-commit verification for DFlash & exact & Bitwise with matched pools; one session: $+$6--12\% at $c=8$--32, $-$3.2\% at $c=1$~\ev{drafter-fold}\\
& FP8 W8A8 weights, FP8 KV; int4 target with a matched drafter~\cite{nota2026} & lossy & FP8 weights: no consistent gain; FP8 KV: no gain (single runs); int4 \pend{moonshot-repeat}\\
Bytes per draft step & Hot-vocabulary or low-precision draft head & exact & \pend{moonshot-patches}\\
Concurrency and state & One state slot per request (lifts the 133-request cap) & exact & \pend{bench-capacity}\\
& FP16 state; replayed decoding; rounding-preserving replay (P4) & lossy, exact up to rounding, exact & FP16: $1.24\times$ at $c=128$ (single run, quality pending); replayed decoding: $1.08\times$ at $c=128$ on the confirmation split in one session~\ev{bench-frontier}, outputs rounding-level~\ev{bench-equality}, more sessions \pend{bench-replayssm}; P4: kernel $1.22\times$ at $B=128$, derived about $1.08\times$ per decode step~\ev{p4-replay}; served test \pend{moonshot-p4}\\
Backbone GEMMs and norms & Kernel routing per projection and $M$ (Hopper GEMV, Triton with programmatic dependent launch, cuBLASLt); packed GDN input projection; norm and SiLU folded into the GEMM & exact up to rounding (packed projection alone: identical outputs at $c=1$) & Best isolated kernel up to $M=16$: 0.77--0.92 of cuBLAS's time on the output projections, 0.91--1.00 on the other large ones; no folding variant pays; served on tuned plain decoding, one session: 3.4\% faster at $c=1$, 0.4\% at 8, 1.0\% at 128~\ev{backbone}; against tuned MTP \pend{backbone-served}\\
System overheads & Host gaps and synchronizations, fusion, attention backend per shape & exact; Triton target attention exact up to rounding & Measured shares in Table~\ref{tab:profile-ranked}; Triton attention: MTP $1.18\times$ at $c=1$, draft only $1.08\times$ (single runs~\ev{host-levers}); blocking planning reads removed with token-identical outputs: tuned MTP $1.075$--$1.096\times$ per-request rate served at $c=1$--32 (interleaved, one session), held cycle 5.9--8.8\% shorter~\ev{hostgap}; DFlash with FlashInfer draft attention \pend{hostgap-served}\\\bottomrule
\end{tabularx}
\caption{Levers by axis. A change of the target's attention backend changes the target's arithmetic, so it is classed by its greedy outputs like any other lever. $^\ast$The replay fold is written as a bitwise clone of the verify kernel's state update, but the verify outputs come from a chunked transform on TF32 tensor cores, whose rounding differs from the recurrent kernel's by more than a reassociation; that mechanism suggested lossy, but every first divergence of the buffered MTP arm from stock MTP is rounding-level~\ev{bench-equality}.}
\label{tab:axes}
\end{table}

\paragraph{What did not work}\label{sec:stack-failed} FP8 W8A8 weights, whose CUTLASS GEMM in the installed aarch64 \code{sgl-kernel} aborts at run time and whose Triton fallback showed no consistent gain in single runs; FP8 KV cache, which saves nothing at short contexts; BF16 recurrent state, which moves the same bytes as FP16 and is worse on quality in published measurements on another GDN model~\cite{damp}; folding the RMSNorm and the SiLU into the next GEMM's prologue, 12--90~\textmu s per layer slower up to $M=32$ and 82--543~\textmu s at 64 and 128 as built, with the single-pass and SiLU-only variants slower at every $M$ from 2 to 16, in isolated layer skeletons~\ev{backbone}; and 2:4 sparsity, which the engine does not support at the pin (its compressed-tensors loader raises an error for 2:4 schemes; code reading, \path{python/sglang/srt/layers/quantization/compressed_tensors/compressed_tensors.py}). An MTP server with BF16 state at 512 concurrent requests also failed to capture its verify graph, on a FlashInfer workspace overflow (Appendix~\paperref{app:serving-protocol} of the paper). Two proposals were rejected by their pre-registered tests: sharing computation across requests (P5, above) and a certified decoder tail (P1, Section~\ref{app:tail}). The repeated runs, with the mechanism of any lever that fails or regresses in them, are \pend{moonshot-repeat}.

\subsection{Interactions and the two stacks}\label{app:stacks}\label{sec:stack-stacks}
Speedups from separate levers do not multiply, and quality losses need not add. Three interactions are predictable from the mechanism. A lossy target lowers acceptance, because the drafter was trained against the BF16 target; the published int4 recipe fine-tunes its drafter against the int4 target for exactly this reason~\cite{nota2026}. A cheaper verifier makes drafting relatively more expensive, which raises the value of a cheap draft head. And every lever that shrinks per-request bytes also raises capacity, so its measured gain depends on whether the server was allowed to use that capacity. Each pair is therefore measured in the baseline, $A$, $B$, $A{+}B$ pattern rather than by multiplying isolated speedups.

The \emph{exact stack} admits only levers that keep the target's decisions up to rounding-level near ties: capacity and scheduling changes, drafting levers, the levers whose greedy outputs are exact up to rounding (replayed verification among them), and the certified head. The \emph{lossy stack} admits everything, under a quality budget declared before measuring: at most 1.0 point of GSM8K accuracy below the reference on the full test split with thinking on (paired, with an exact McNemar test), teacher-forced top-1 agreement of at least 98\% and mean top-20 KL divergence of at most 0.01 nats on a fixed probe set of 48 prompts and 256 tokens, each reported next to the reference's own run-to-run noise. The combination's quality is measured on the combined configuration, never summed from its parts. Proposition~\ref{prop:kl} explains why the budget is empirical: a per-token logit error small enough to certify a long output as close in total variation is far below what any low-precision lever delivers. The interaction matrix (speed at concurrency 1 and at peak, acceptance and quality for each measured pair) and both frontiers, from one server launch per arm, are \pend{moonshot-frontiers}.

\subsection{A sensitivity workload for output length}\label{notes:sensitivity}
\paragraph{Sensitivity to output length} With a closed-loop client and equal output lengths, plain requests start and finish in synchronized waves, while speculative arms advance requests at rates that vary with acceptance, so their completions spread out into more, smaller prefill passes and a drain at the end of a point. A sensitivity workload, declared before any arm was timed on it, removes the equal lengths and nothing else: each confirmation prompt is given its own greedy completion length under plain decoding with natural stopping, capped at 2,048 tokens, so that every arm generates the same number of tokens per request. It runs at concurrency 32 and 128 with $\max(64,8c)$ measured requests, the ordered prefix of the confirmation split (256 and 1,024 prompts). The arms follow a rule fixed before any sensitivity run: for each speculative family and concurrency, the arm with the highest mean $y$ over the three confirmation sessions, invalid points excluded and only arms with a valid point in all three eligible, together with the family's runner-up if it is within 2\% of that mean. Each runs with its matched plain baseline, in three sessions. The fixed-length panel stays the primary result; where the sensitivity workload changes a ranking, both numbers are reported. Its results are \pend{bench-sensitivity}.
